\chapter{Introdução}\label{cap:introducao}
% ===========================================================================

Este documento apresenta o projeto de formatura intitulado ``Aplicação de
Protocolos de Oblivious Transfer para Preservação de Privacidade em Transações
Digitais''.

\section{Motivação}

No cenário atual do comércio eletrônico, a coleta massiva de dados se tornou
uma prática recorrente, de forma que qualquer compra torna a navegação pela
internet uma experiência infestada de anúncios de produtos relacionados.
Embora a vantagem de marketing dessa prática seja inegável, fornecer para os
consumidores a privacidade que a maioria das lojas virtuais atualmente viola é
um aspecto de grande interesse e necessidade para a sociedade.

A literatura de economia da privacidade descreve como as informações pessoais
geradas em transações \emph{online} passaram a ter valor comercial próprio e
como o consumidor perde o controle sobre quem as usa e com que finalidade
\cite{acquisti2016}. No Brasil, a Lei Geral de Proteção de Dados Pessoais
disciplina o tratamento desses dados \cite{lgpd2018}, mas a própria execução de
uma compra é base legal para esse tratamento: a loja continua sabendo o que
cada cliente comprou.

A criptografia oferece ferramentas para que a transação aconteça sem que essa
informação chegue ao vendedor. As assinaturas cegas, propostas por
\citeonline{chaum1983} para pagamentos eletrônicos não rastreáveis, permitem
que uma parte assine um valor sem conhecer seu conteúdo. Os protocolos de
transferência inconsciente (\emph{Oblivious Transfer} -- OT), introduzidos por
\citeonline{rabin2005}, permitem que um receptor obtenha um entre vários itens
sem que o emissor saiba qual. \citeonline{aiello2001} aplicaram o OT
diretamente à venda de bens digitais, com o OT com preço (\emph{priced
oblivious transfer}), em que o comprador paga e recebe itens sem que o
vendedor saiba quais foram escolhidos. Trabalhos posteriores estenderam a
ideia para escolhas adaptativas, feitas uma após a outra \cite{camenisch2007},
e para bases de dados em que cada registro tem a sua própria política de
acesso \cite{camenisch2009}.

\section{Objetivos}

\textbf{Objetivo Geral:} Desenvolver um protótipo funcional de um sistema de
venda de ativos digitais em formato de catálogo, utilizando criptografia
baseada em \emph{Oblivious Transfer} (OT, Transferência Inconsciente) para passar chaves
de criptografia simétrica usadas para descriptografar os bens digitais
adquiridos de Rede de Distribuição de Conteúdo (Content Delivery Network --
CDN), garantindo que o servidor da loja seja incapaz de identificar quais
produtos foram acessados ou adquiridos por um comprador específico.

\textbf{Objetivos Específicos:}

\begin{itemize}
\item Modelar a arquitetura do sistema para permitir a transação segura onde a
escolha do cliente é ofuscada para o servidor.
\item Implementar o algoritmo criptográfico escolhido (como 1-out-of-N) em uma
solução com 3 componentes: site online, rede de distribuição de conteúdo (CDN)
e programa local instalado no computador do cliente.
\item Desenvolver mecanismos para validar a legitimidade da transação sem
revelar o conteúdo da compra.
\item Realizar testes de segurança e desempenho para validar a impossibilidade
de rastreamento de preferências.
\end{itemize}

\section{Justificativa}

O projeto propõe uma solução técnica para um problema crescente de privacidade
digital. O sistema permite que o comprador adquira produtos digitais dentre
várias opções, tendo acesso a apenas uma, mas de forma que o vendedor
permaneça matematicamente incapaz de saber qual item específico foi
selecionado pelo cliente.

A solução implementa algoritmos de Oblivious Transfer
integrados a uma interface de venda, garantindo que a loja receba o pagamento,
mas não consiga registrar históricos de compras, preferências ou interesses do
usuário, assegurando assim um elevado nível de privacidade para os
consumidores.

Os trabalhos de \citeonline{aiello2001}, \citeonline{camenisch2007} e
\citeonline{camenisch2009} concentram-se nos protocolos e em suas provas de
segurança. Este trabalho integra o OT a um fluxo de compra completo, com
pagamento por \emph{gateway}, distribuição dos livros cifrados por uma CDN e
remuneração da CDN por meio de tokens com assinatura cega \cite{chaum1983}, e
avalia o resultado contra requisitos de desempenho e de privacidade. Adota,
ainda, uma simplificação de engenharia: em vez de esconder o preço de cada
item, como no OT com preço, agrupa os livros por preço e cobra pelo grupo e
pela quantidade de itens.

\section{Organização do Trabalho}

Este documento está estruturado da seguinte forma: o \autoref{cap:conceitos}
apresenta os conceitos fundamentais que embasam o projeto. O
\autoref{cap:metodo} descreve o método do trabalho e suas fases. O
\autoref{cap:especificacao} define os requisitos do sistema, o modelo de
ameaça e os casos de uso. O \autoref{cap:desenvolvimento} apresenta as
tecnologias utilizadas, o projeto e a implementação da solução e os testes
realizados. O \autoref{cap:consideracoes} apresenta as conclusões, as
contribuições e as perspectivas de continuidade.
